%% ++++++++++++++++++++++++++++++++++++++++++++++++++++++++++++
%% Hauptdatei, Wurzel des Dokuments
%% ++++++++++++++++++++++++++++++++++++++++++++++++++++++++++++

% Headerfeld, Typ des Dokumentes, einzubindende Packages.
% Hier bei Bedarf Änderungen vornehmen.
\documentclass
[   twoside=false,     % Einseitiger oder zweiseitiger Druck?
    fontsize=12pt,     % Bezug: 12-Punkt Schriftgröße
    DIV=15,            % Randaufteilung, siehe Dokumentation "KOMA"-Script
    BCOR=17mm,         % Bindekorrektur: Innen 17mm Platz lassen. Copyshop-getestet.
%    headsepline,
    headsepline,  % Unter Kopfzeile Trennlinie (aus: headnosepline)
    footsepline,  % Über Fußzeile Trennlinie (aus: footnosepline)
    open=right,        % Neue Kapitel im zweiseitigen Druck rechts beginnen lassen
    paper=a4,          % Seitenformat A4
    abstract=true,     % Abstract einbinden
    listof=totoc,      % Div. Verzeichnisse ins Inhaltsverzeichnis aufnehmen
    bibliography=totoc,% Literaturverzeichnis ins Inhaltsverzeichnis aufnehmen
    titlepage,         % Titelseite aktivieren
    headinclude=true,  % Seiten-Head in die Satzspiegelberechnung mit einbeziehen
    footinclude=false, % Seiten-Foot nicht in die Satzspiegelberechnung mit einbeziehen
    numbers=noenddot   % Gliederungsnummern ohne abschließenden Punkt darstellen
]   {scrreprt}         % Dokumentenstil: "Report" aus dem KOMA-Skript-Paket

\usepackage[active]{srcltx}
%\usepackage[activate=normal]{pdfcprot} % Optischer Randausgleich -> pdflatex!
\usepackage{ifthen}
\usepackage[ngerman]{babel}   % Neue Deutsche Rechtschreibung
%\usepackage[latin1]{inputenc} % Zeichencodierung nach ISO-8859-1
\usepackage[utf8]{inputenc}   %	Zeichencodierung nach UTF-8 (Unicode)
\usepackage[T1]{fontenc}
%\usepackage{ae} % obsolet und durch lmodern ersetzt
\usepackage{lmodern}
\usepackage[T1]{url}
\usepackage[final]{graphicx}
\RequirePackage{scrlfile}
\usepackage{skak}
\usepackage{amssymb}
\ReplacePackage{scrpage2}{scrlayer-scrpage}
% old: \usepackage[automark]{scrpage2}
\usepackage[automark]{scrlayer-scrpage}
\usepackage{setspace}
%\usepackage[first,light]{draftcopy} % Für Probedruck
\usepackage[plainpages=false,pdfpagelabels,hypertexnames=false]{hyperref}

% Tiefe der Kapitelnummerierung beeinflussen
\setcounter{secnumdepth}{3} % Tiefe der Nummerierung
\setcounter{tocdepth}{3}    % Tiefe des Inhaltsverzeichnisses

% Datum anpassen
\newcommand{\leadingzero}[1]{\ifnum #1<10 0\the#1\else\the#1\fi}    
\renewcommand{\today}{\leadingzero{\day}.\leadingzero{\month}.\the\year}     % DD.MM.YYYY

% Hier in die zweite geschweifte Klammer jeweils
% die persönlichen Daten und das Thema der Arbeit eintragen:
\newcommand{\namedesautors}{Jonathan Skopp}
\newcommand{\artderausarbeitung}{\namedesautors}
\newcommand{\themaderarbeit}{Endspieltraining}

% FEN
\newcommand{\white}[1]{
	\centering
	\newgame
	\fenboard{#1} 
	\showboard \hspace{1cm} $\square$ \\
	FEN: #1 
}

\newcommand{\black}[1]{
	\centering
	\newgame
	\fenboard{#1} 
	\showboard \hspace{1cm}  $\blacksquare$\\
	FEN: #1 
}

% Faulheit
\newcommand{\Remis}{$\frac{1}{2}-\frac{1}{2}$}
\newcommand{\unterstreichen}[1]{\underline{\underline{#1}}}
\newcommand{\Kommentar}[1]{\glqq #1\grqq{}}
\newcommand{\links}{\raggedright}

% PDF Metadaten definieren
\hypersetup{
   pdftitle={\themaderarbeit},
   pdfsubject={\artderausarbeitung},
   pdfauthor={\namedesautors},
   pdfkeywords={\artderausarbeitung; TU-Ilmenau}}


% Abkürzungsverzeichnis beeinflussen. Hier nichts ändern!
\usepackage[intoc]{nomencl}
  \AtBeginDocument{\setlength{\nomlabelwidth}{.25\columnwidth}}
  \let\abbrev\nomenclature
  \renewcommand{\nomname}{Abkürzungsverzeichnis und Formelzeichen}
  \renewcommand{\nomlabel}[1]{#1 \dotfill}
  \setlength{\nomitemsep}{-\parsep}
  \makenomenclature

\usepackage[normalem]{ulem}
  \newcommand{\markup}[1]{\textbf{#1}}

% Seitenlayout festlegen. Hier nichts ändern!
\pagestyle{scrplain}
\ihead[]{\headmark}
\ohead[]{\pagemark}
\chead[]{}
\ifoot[]{}
\ofoot[]{\scriptsize  \namedesautors}
\cfoot[]{}
\renewcommand{\titlepagestyle}{scrheadings}
\renewcommand{\partpagestyle}{scrheadings}
\renewcommand{\chapterpagestyle}{scrheadings}
\renewcommand{\indexpagestyle}{scrheadings}

% Abschnittsweise Nummerierung anstatt fortlaufend. Hier nichts ändern!
\makeatletter
\@addtoreset{equation}{chapter}
\@addtoreset{figure}{chapter}
\@addtoreset{table}{chapter}
\renewcommand\theequation{\thechapter.\@arabic\c@equation}
\renewcommand\thefigure{\thechapter.\@arabic\c@figure}
\renewcommand\thetable{\thechapter.\@arabic\c@table}
\makeatother

% Quelltextrahmen, klein. Hier nichts ändern!
\newsavebox{\inhaltkl}
\def\rahmenkl{\sbox{\inhaltkl}\bgroup\small\renewcommand{\baselinestretch}{1}\vbox\bgroup\hsize\textwidth}
\def\endrahmenkl{\par\vskip-\lastskip\egroup\egroup\fboxsep3mm%
\framebox[\textwidth][l]{\usebox{\inhaltkl}}}

% Quelltextrahmen, normale Groesse. Hier nichts ändern!
\newsavebox{\inhalt}
\def\rahmen{\sbox{\inhalt}\bgroup\renewcommand{\baselinestretch}{1}\vbox\bgroup\hsize\textwidth}
\def\endrahmen{\par\vskip-\lastskip\egroup\egroup\fboxsep3mm%
\framebox[\textwidth][l]{\usebox{\inhalt}}}

% Andere Befehle.
\newcommand{\entspricht}{\stackrel{\wedge}{=}}

% Tabellenspaltendefinitionen mit fester Breite --> somit Zeilenumbruch innerhalb einer Zelle möglich
% aus http://www.torsten-schuetze.de/tex/tabsatz-2004.pdf
\usepackage{array, booktabs}
\newcolumntype{f}{>{$}l<{$}}
\newcolumntype{n}{>{\raggedright}l}
\newcolumntype{N}{>{\scriptsize}l}
\newcolumntype{v}[1]{>{\raggedright\hspace{0pt}}m{#1}}
\newcolumntype{V}[1]{>{\scriptsize\raggedright\hspace{0pt}}m{#1}}
\newcolumntype{Z}[1]{>{\raggedright\centering}m{#1}}
\newcolumntype{k}[1]{>{\raggedright}p{#1}}
% ergibt Tabllenspalte fester Breite, linksbündig
% Umbruch innerhalb der Zelle mit \\, neue Tabellezeile mit \tabularnewline
% \addlinespace für Gruppentrennung (aus \texttt{booktabs.sty})


\begin{document}
\onehalfspacing
\begin{titlepage}
\centering
\includegraphics[scale=0.5]{DSB.png}\\[3ex]
{\Large \textsc{Deutscher Schachbund}}\\[3ex]
{\Large Thüringer Schachbund}\\[3ex]
\vfill
{\Large \textbf{\artderausarbeitung}}\\[4ex]
{\large \textbf{\themaderarbeit}}\\[5ex]
\vfill
\begin{tabular}{rl}
\hline \\
vorgelegt von:          & \quad \namedesautors\\[1,5ex]
eingereicht am:         & \quad 10.\,11.\,2020\\[1,5ex]
geboren am:             & \quad 28.\,04.\,1998 in Sondershausen\\[1,5ex]
Ausbildung:             & \quad C-Trainer Lehrgang 2020\\[1,5ex]
Hausarbeit:             & \quad Lehrprobe\\[1,5ex]
                        & \quad Starke und schwache Felder\\[1,5ex]
Dozent:                 & \quad Norbert Reichel (B-Trainer)\\[1,5ex]

\end{tabular}
\vfill
\end{titlepage}








%% ++++++++++++++++++++++++++++++++++++++++++++++++++++++++++++
%% Zusammenfassung, Abstract
%% ++++++++++++++++++++++++++++++++++++++++++++++++++++++++++++


\renewcommand{\abstractname}{Kurzfassung}
\begin{abstract}
\ldots \emph{Hier später die eigene deutsche Kurzfassung einfügen}\ldots
\par{}
Dieses Dokument soll als Gerüst für eigene \LaTeX\ Dokumente dienen und
gleichzeitig Beispiele für häufig verwendete Konstrukte wie Tabellen,
Formeln oder Grafiken liefern. Es empfielt sich, diese Elemente per
Cut\&Paste zu kopieren und einzufügen.
\end{abstract}




% Inhaltsverzeichnis
\cleardoublepage % Seitenumbruch erzwingen vor Änderung des Nummerierungsstils
\pagenumbering{arabic} % Nummerierung der Seiten ab hier: i, ii, iii, iv...
\pagestyle{scrheadings} % Ab hier mit Kopf- und Fusszeile
\tableofcontents


% Training
\chapter{Training 08.12.2020}
 Die folgenden Aufgaben beziehen sich auf die erste Lichess Studie \cite{Study1}.
 
\section{Kapitel 1}
\subsection{Mit Schwarz}

\black{8/3k4/8/3K4/3P4/8/8/8 b - - 0 1}


\links
Der weiße Bauer auf dem Feld b4 hat die Mittellinie noch nicht überschritten. \\
Die Schlüsselfelder des weißen Bauerns auf b4 sind c6, d6 und e6.\\
 Wenn der weiße König diese erreicht, ist die Partie für weiß gewonnen. Schwarz muss deshalb dies verhindern. Dies ist ihm aber nicht möglich. \\
 Sobald der König auf e7 zieht deckt er das Schlüsselfeld c6 nicht mehr. \\  Folglich kann schwarz nicht verhindern, dass Weiß eines der begehrten Felder erreicht. Demnach verliert er die Partie

\pagebreak

\subsection{Mit Weiß}
\white{8/4k3/8/3K4/3P4/8/8/8 w - - 0 2}


\links
Nach Ke7 gibt dieser das Schlüsselfeld c6 frei. Dieses kann Weiß nun besetzen. Dadurch hat er das Schlüsselfeld erreicht und gewinnt die Partie. \\

\centering
\begin{tabular}[h]{c|c}
	\textbf{Weiß}  & \textbf{Schwarz} \\
	\hline
	Kc6 & Kd8 \\
	d5 & \unterstreichen{1-0}
\end{tabular}


\section{Kapitel 3}

\white{k7/8/2P5/8/K7/8/8/8 w - - 0 1}


\links
Weiß möchte in dieser Partie auf Gewinn spielen. \\
Der Bauer hat die Mittellinie bereits überschritten, daher sind die Schlüsselfelder die Matrix von b7 bis d8. Um diese zu erreichen, spielt Weiß hier Kb5, denn er möchte um den Bauern herumlaufen. Die Opposition braucht weiß nicht zu fürchten, da der weiße Bauer das Feld b7 deckt. Nach Kb8 kann Weiß mit dem König nach b6 ziehen. Um an dem Bauern zu bleiben zieht Schwarz Kc8. Daraufhin kann ich c7 spielen und die Partie mit einem Punkt für mich beenden. \\
Wenn der Bauer den König mit Schach zwingt unter ihn zu ziehen, endet die Partie unentschieden. Zwingt der Bauer den König ohne Schach unter den Bauern, endet die Partie gewonnen. So auch hier. \\

\centering
\begin{tabular}[h]{c|c}
	\textbf{Weiß}  & \textbf{Schwarz} \\
	\hline
	Kb5 & Kb8\\
	c7 & \unterstreichen{1-0}
\end{tabular}


\section{Kapitel 4}
\black{8/k7/p1K5/8/1P6/8/8/8 b - - 0 1}


\links
Hier ist schwarz am Zug.\\
Dieser ist offensichtlich mit einem Remis glücklich. Als mögliche Züge kommt hier nur a5 in Frage. Nun hat Weiß 2 Optionen. Entweder er schlägt diesen oder er ignoriert ihn. Wenn er ihn schlägt, ist es ein Randbauer und der schwarze König hat das Eckfeld. Also spielt Weiß b5.\\
Schwarz möchte nun den Bauern auf der B-Linie blockieren. Dazu muss er diese betreten. Dies geht nur auf dem Feld b8. Wenn er dies nicht jetzt tut, kann er es gar nicht mehr tun, denn Weiß riegelt das Feld sonst mit Kc7 ab.
Daher zieht schwarz Kb8. \\
Weiß muss sich jetzt um dem a Bauern kümmern und geht über c5-b4 an den laufenden Bauern heran und schlägt diesen. Der weiße König steht nun auf a3.\\ 
schwarz kann nun ebenfalls an den Bauern herranlaufen. Nun stehen sich die Beiden Könige in einer Opposition gegenüber. Daher endet diese Partie unentschieden. \\
\centering
\begin{tabular}[h]{c|c}
	\textbf{Weiß}  & \textbf{Schwarz} \\
	\hline
	$\cdots$ & a5\\
	b5 & Kb8 \\
	Kc5 & a4 \\
	Kb4 & a3 \\
	Kxa3 & Kb7 \\
	Kb4 & Kb6 \\
	\unterstreichen{\Remis}
\end{tabular}

\section{Kapitel 5}
\white{6k1/8/3K2p1/8/6P1/8/8/8 w - - 0 1}\\
\links
In dieser Stellung ist zu erkennen, dass der Bauer auf g4 markant steht. Nach dem Zug g5 sperrt dieser Bauer dem schwarzen König die wichtigen Fluchtfelder h6 und f6 ab. Der weiße König bedrängt ihn zusätzlich von der d bzw. f Linie. Dadurch muss der schwarze  König langfristig den Bauern auf g6 aufgeben.\\
\centering
\begin{tabular}[h]{c|c}
	\textbf{Weiß}  & \textbf{Schwarz} \\
	\hline
	g5 & Kf7 \\
	Kd7 & Kg7 \\
	Ke7 & Kh7 \\
	Kf7 & Kh8 \\
	Kxg6 & Kg8 \\
	Kh6 & Kh8 \\
	g6 & Kg8 \\
	g7 & $\cdots$\\
	\unterstreichen{1-0}
\end{tabular}

\section{Kapitel 6}
\white{8/ppp5/8/PPP5/8/7k/8/7K b - - 0 1}\\
\links

Hier wird recht schnell deutlich, dass man einen Durchbruch erzeugen muss. Daher weis ich, dass ich mit b6 anfangen muss. Wenn nun einer der beiden benachbarten Bauern von schwarz den b6 schlägt, kann ich mir einen Freibauern bilden. Dieser läuft dann durch und kann sich umwandeln. \\
\centering
\begin{tabular}[h]{c|c}
	\textbf{Weiß}  & \textbf{Schwarz} \\
	\hline
	b6 & cxb6\\
	a6 & bxa6\\
	c6 & a5 \\
	c7 & Kh4\\
	c8D & $\cdots$ \\
	\unterstreichen{1-0}
\end{tabular}
\clearpage

\section{Kapitel 7}
\black{r3r1k1/5ppp/q7/8/1B2n3/R1P4Q/5PPP/5RK1 b - - 0 1}\\
\links

In dieser Matt-Aufgabe fällt auf das schwarz viele aktive Figuren am Königsflügel hat. Ferner fällt der Turm auf e8 auf, welcher die gesamte e-Linie komplett kontrolliert.  Nun gilt es nur noch den Schlag zu finden.\\
\centering

\begin{tabular}[h]{c|c}
	\textbf{Weiß}  & \textbf{Schwarz} \\
	\hline
	$\cdots$ & Dxf1+ \\
	Kxf1 & Sd2+ \\
	Kg1 & Ke1\#\\
	$\cdot$ & \unterstreichen{0-1}
\end{tabular}

\section{Kapitel 8}
\white{3r3k/pbq1b1p1/1p3n1p/5Q1P/3n3N/P7/BP3PP1/2B1R2K w - - 0 1}\\
\links

Hier brauche ich nicht viel sagen. Der Springer lacht doch schon. Mit dem Läufer auf a2 schreit das schon nach einem erstickten Matt. \\
\centering
\begin{tabular}[h]{c|c}
	\textbf{Weiß}  & \textbf{Schwarz} \\
	\hline
	Sg6+ & Kh7\\
	Sf8+ & Kh8 \\
	Dh7+ & Sxh7\\
	Sg6\# & $\cdot$\\
	\unterstreichen{1-0}
\end{tabular}

\section{Kapitel 9}
\black{k3r2r/pp4bp/6p1/5p2/N2P4/8/1P3PPP/R1R3K1 w - - 0 1}\\
\links

Diese Aufgabe ist auch recht bekannt. Vor allem wenn man sich die a-Linie betrachtet. Auch hier scheint der Springer schon ausdrücklich nach b6 zu wollen. Der Rest ergibt sich dann.\\
\centering

\begin{tabular}[h]{c|c}
	\textbf{Weiß}  & \textbf{Schwarz} \\
	\hline
	Sb6+ & Kb8 \\
	Sd7+ & Ka8 \\
	Txa7+ & Kxa7 \\
	Ta1\# & $\cdot$\\
	\unterstreichen{1-0}
\end{tabular}

\section{Kapitel 10}
\white{8/p2p4/8/8/8/k7/5P1P/7K w - - 0 1}\\
\links

Weiß kann nicht blind mit dem f oder h Bauern loslaufen. Das liegt daran, dass der König den f Bauern einholen würde und mit d5-d1 genauso schnell ist wie weiß. \\ 
Allerdings kann der König nicht den h Bauern einholen. Interessant ist auch, dass in aktueller Stellung der f-Bauer mit Schach einzieht.\\
Wir spielen nun f4 und immer wenn der König in Richtung der Bauern zieht, ziehen wir den H Bauern. Dadurch können wir unseren f-Bauern mit Schach umwandeln. Der schwarze König schlägt diesen und steht nun ebenfalls auf der achten Reihe. Dadurch können wir mit h8D+ uns eine neue Dame holen und den d-Bauern daran hindern sich umzuwandeln. \\
\centering
\begin{tabular}[h]{c|c}
	\textbf{Weiß}  & \textbf{Schwarz} \\
	\hline
	f4 & Kb4 \\
	h4 & d5 \\
	f5 & Kc5 \\
	h5 & d4 \\
	f6 & Kd6 \\
	h6 & d3 \\
	f7 & Ke7 \\
	h7 & d2 \\
	f8D+ & Kxf8 \\
	h8D+ & $\cdot$\\
	\unterstreichen{1-0}
\end{tabular}

\section{Kapitel 11}
\white{1r3k2/2q2b1p/5Qp1/1p2n3/p3P3/8/PPP3PP/1K1R1R2 w - - 0 1}\\
\links

In dieser Aufgabe ist von vorne nicht zu erkennen worum es eigentlich geht. \\ 
Das Matt mit Dxf7 wird leider von dem Springer und der Dame verhindert. Folglich sollte man schauen ob man diese Ablenken kann.  Nach Td7 muss Dxd7 kommen (Sxd7 $\rightarrow$ Dxf7\#). Dieser Zug lenkt aber die Dame von einem anderen Feld ab; und zwar von dem Turm auf b8. Dieser ist nun ungedeckt und kann mit Dh8+ angegriffen werden. Den Rest erkennt man in der Partie.\\
\centering
\begin{tabular}[h]{c|c}
	\textbf{Weiß}  & \textbf{Schwarz} \\
	\hline
	Td7 & Dxd7 \\
	Dh8+ & Ke7 \\
	Dxe5+ & De6 \\
	Dc7+ & Dd7 \\
	Dxb8 & $\cdot$\\
	\unterstreichen{1-0}
\end{tabular}


\chapter{Training 15.12.2020}
 Die folgenden Aufgaben beziehen sich auf die zweite Lichess Studie \cite{Study2}

\section{Kapitel 1}
\links
Diese Aufgabe ist eine Wiederholung. \\

\subsection{Mit Schwarz}
\black{8/3k4/8/3K4/3P4/8/8/8 b - - 0 1}\\

\links
Wenn in dieser Stellung schwarz am Zug ist, gewinnt Weiß. Das liegt daran, dass er nach Ke7 das Schlüsselfeld c6 freigibt und nach Kc7 das Schlüsselfeld e6. \\
\centering
\begin{tabular}[h]{c|c}
	\textbf{Weiß}  & \textbf{Schwarz} \\
	\hline
	$\cdots$ & Ke7 \\
	Kc6 & Kd8 \\
	Kd6 & Kc8 \\
	Ke7 & kc7 \\
	d5 & $\cdot$ \\
	\unterstreichen{1-0}
\end{tabular}

\subsection{Mit Weiß }
\white{8/3k4/8/3K4/3P4/8/8/8 b - - 0 1}

Wenn Weiß in dieser Stellung am Zug ist, endet die Partie \unterstreichen{\texttt{Remis}}, da schwarz immer alle Schlüsselfelder abdecken kann.

\section{Kapitel 2}
\black{8/8/8/6p1/7k/8/6K1/8 b - - 0 1}\\
\links

In dieser Stellung gibt es nur 2 Züge die interessant sind:
 \begin{itemize}
	\item Kg4 
	\item g4
\end{itemize}
Nach g4 ist die Partie leider nach der Definition aus 1.1.2 Remis. Deswegen muss hier Kg4 kommen und dann ist das eine Frage des Zugzwangs. Denn nach Kg4 bekomme ich immer ein Schlüsselfeld.\\
\centering
\begin{tabular}[h]{c|c}
	\textbf{Weiß}  & \textbf{Schwarz} \\
	\hline
	$\cdots$ & Kg4 \\
	Kh2 & Kf3 \\
	Kh3 & g4+ \\
	Kh2 & g3+ \\
	Kh1 & Kf2 \\
	$\cdot$ & \unterstreichen{0-1}
\end{tabular}	

\section{Kapitel 3}
\black{8/8/8/8/8/7p/7k/5K2 b - - 0 1}\\
\links
Diese Stellung ist \unterstreichen{\texttt{Remis}}, da Weiß immer das Feld h2 decken kann. Schwarz kann dagegen nichts machen. \\
\centering
\begin{tabular}[h]{c|c}
	\textbf{Weiß}  & \textbf{Schwarz} \\
	\hline
	$\cdots$ & Kg3 \\
	Kg1 & h2+\\
	Kh1 & Kf2 \\
	$\cdot$ & \unterstreichen{\Remis}
\end{tabular}	

\section{Kapitel 4}
\white{8/2k5/3p4/1K1P4/8/8/8/8 w - - 0 1}\\
\links
Hier ist die Idee, dass ich mit dem König von \Kommentar{links} an den König herranlaufe und diesen so abdränge. Die beiden Widderbauern decken jeweils die Felder neben den des anderen Bauerns. Dies ist der Witz der Aufgabe, da ich so den König von dem Bauern abdrängen kann. \\
\centering
\begin{tabular}[h]{c|c}
	\textbf{Weiß}  & \textbf{Schwarz} \\
	\hline
	Ka6 & Kd7 \\
	Kb7 & Ke7 \\
	Kc7 & Ke8 \\
	Kxd6 & Kd8 \\
	Ke6 & Ke8 \\
	d6 & Kd8 \\
	d7 & $\cdot$ \\
	\unterstreichen{1-0}
\end{tabular}

\section{Kapitel 5}
\black{8/8/8/8/1P6/6k1/8/6K1 b - - 0 1}\\
\links
 Die schöne \Kommentar{Quadratregel}. \\
 Das Quadrat des Bauerns geht von b4-b8-f4-f8. Da schwarz dran ist, kann er dieses Quadrat betreten und so den Bauern in der Umwandlung schlagen.\\
 \centering
 \begin{tabular}[h]{c|c}
 	\textbf{Weiß}  & \textbf{Schwarz} \\
 	\hline
 	$\cdots$ & Kf4 \\
 	b5 & Ke5 \\
 	b6 & Kd6 \\
 	b7 & Kc7 \\
 	\unterstreichen{\Remis} & \unterstreichen{\Remis}
 \end{tabular}

\section{Kapitel 6}
\white{8/5p2/4p3/8/3P4/5k2/P7/5K2 w - - 0 1}\\
\links
Hier geht das Quadrat des a2 Bauerns über folgende Felder:  a2-g2-a8-g8. \\
Dadurch würde der Schwarze König den Bauern gewinnen können. Durch den Zug g5 stelle ich einen schwarzen Bauern auf das Feld. \\
\centering
\begin{tabular}[h]{c|c}
	\textbf{Weiß}  & \textbf{Schwarz} \\
	\hline
	d5 & exd5 \\ 
	a4 & Ke4 \\
	a5 & $\cdot$ \\
	\unterstreichen{1-0}
\end{tabular}

\section{Kapitel 7}
\white{8/8/8/3kp3/8/8/8/3K4 w - - 0 1}\\
\links

Der Bauer hat die Mitte bereits übertreten, dadurch sind die Schlüsselfelder d3-d2-f2-f3.\\
Weiß muss erreichen, das er die Oposition erzwingt. Sollte er direkt auf die zweite Reihe gehen, ist dies nicht möglich. Daher zieht er erst Ke1. Nun muss schwarz dies erlauben.\\

\centering
\begin{tabular}[h]{c|c}
	\textbf{Weiß}  & \textbf{Schwarz} \\
	\hline
	Ke1 & Ke4 \\
	Ke2 & $\cdot$ \\
	\unterstreichen{1-0}
\end{tabular}

\section{Kapitel 8}
\white{5k2/8/8/8/1P6/8/8/3K4 w - - 0 1}\\
\links
Weiß muss offensichtlich an die Bauern heranlaufen. Dies ist auf 2 Feldern möglich. c4 und a4. Für das Feld benötigt er 3 Halbzüge und schwarz mehr. Das bedeutet er ist schneller da.\\
Dies stellt den Zeitvorteil dar. \\
\centering
\begin{tabular}[h]{c|c}
	\textbf{Weiß}  & \textbf{Schwarz} \\
	\hline
	Kc2 & Ke8 \\
	Kb3 & Kd8 \\
	Ka4 & Kc8 \\
	Ka5 & Kb8 \\
	Kb6 & $\cdot$ \\
	\unterstreichen{1-0}
\end{tabular}

\section{Kapitel 10}
\white{k7/8/1p6/8/1P6/5K2/8/8 w - - 0 1}\\
\links

Die Aufgabe ist eine Wiederholung der Aufgabe 1.5 . Nach b5 kann ich den König wieder erfolgreich abdrängen.  \\
\centering
\begin{tabular}[h]{c|c}
	\textbf{Weiß}  & \textbf{Schwarz} \\
	\hline
	b5 & Kb7 \\
	Ke3 & $\cdot$ \\
	\unterstreichen{1-0}
\end{tabular}

\section{Kapitel 11}
\white{8/7p/8/4kpPP/8/8/8/5K2 w - - 0 1}\\
\links

Die Idee ist, dass schwarz nicht mit dem König von unten an die Bauern herankommt, da diese sonst durchlaufen. Wenn ich mit dem König ziehe, kommt sofort f4 und ich kann meine Bauern nicht mehr verteidigen (Remis). \\
Wenn ich direkt mit g6 loslaufe holt mich der König dennoch ein. Folglich bleibt mir nur der Zug h6 mit dem Opfer auf g6. Grundlegend ist es das selbe, bloß das der h Bauern nun bereits ein Feld weiter vorn steht und sich so umwandeln kann. 

\centering
\begin{tabular}[h]{c|c}
	\textbf{Weiß}  & \textbf{Schwarz} \\
	\hline
	h6 & Ke6 \\
	g6 & $\cdot$ \\
	\unterstreichen{1-0}
\end{tabular}

\section{Kapitel 12}
\white{8/8/3p4/k7/P1P5/p7/2K5/8 w - - 0 1}\\
\links

Offensichtlich steht weiß hier etwas seltsam. Der Bauer auf a3 ist fast durch und die eigenen Bauern stehen irgendwie fremdartig.  Dennoch ist der erste Zug recht offensichtlich. \\
Kb3. Wenn schwarz nun b2 spielt, schlage ich diesen nicht raus sondern spiele Kb2. Das liegt daran, dass ich dann kein Zugzwang wegen dem c Bauer habe. Erst wenn schwarz jetzt den Bauern umwandelt tausche ich diesen. 
Daraufhin nimmt schwarz den a4 und gibt mir so die Zeit c5 zu spielen um die Partie Remis zu halten. \\
Denn nach c5 dxc5 ist es wieder eine definierte Stellung.\\
\centering

\begin{tabular}[h]{c|c}
	\textbf{Weiß}  & \textbf{Schwarz} \\
	\hline
	Kb3 & a2 \\
	Kb2 & a1D+ \\
	Kxa1 & Kxa4 \\
	c5 & dxc5\\
	Ka2 & $\cdot$ \\
	\unterstreichen{\Remis} & \unterstreichen{\Remis}
\end{tabular}

\section{Kaptiel 13}

\white{8/8/3p4/2k5/4P3/8/8/1K6 w - - 0 1}

Der erste Zug e5 stammt aus vorherigen Partien. Nun geht es darum nicht in eine Opposition zu gelangen. Dies ist nur auf dem Feld Kc1 möglich. Sollte schwarz hier die Opposition aufstellen verliert Weiß, sonst ist es Remis. \\
\centering
\begin{tabular}[h]{c|c}
	\textbf{Weiß}  & \textbf{Schwarz} \\
	\hline
	e5 & dxe5\\
	Kc1 & Kc4 \\ 
	Kc2 & $\cdot$ \\
	\unterstreichen{\Remis} & \unterstreichen{\Remis}
\end{tabular}



\chapter{Training 22.12.2020}
 Die folgenden Aufgaben beziehen sich auf die zweite Lichess Studie \cite{Study2}
 
 
\section{Kapitel 14}
\white{8/1k6/1p6/1K6/2P5/8/P7/8 w - - 0 1}\\
\links

Man könnte überlegen, mit a4 loszulaufen nur um dann festzustellen, dass der König jedes mal im Weg steht. \\
Die Idee ist also mit a3 dem schwarzen ein Tempo zu geben um so das Feld a6 zu bekommen. Dann können beide Partein um den Bauern herumlaufen und der c4 Bauer sorgt dann dafür, dass schwarz den Bauern aufgeben muss.\\
 \centering
\begin{tabular}[h]{c|c}
	\textbf{Weiß}  & \textbf{Schwarz} \\
	\hline
	a3 & Kc7 \\
	Ka6 & Kc6 \\
	a4 & Kc7 \\
	Ka7 & Kc6 \\
	Kb8 & Kd6 \\
	Kb7 & Kc5 \\
	Kc7 & Kxc4 \\
	Kxb6 & $\cdot$ \\
	\unterstreichen{1-0}
\end{tabular}

\section{Kapitel 15}
\white{4k3/8/8/1p5p/1P5P/8/8/4K3 w - - 0 1}\\
\links
Beide Könige laufen sofort auf die vierte bzw sechste Reihe los. Dort stehen sie sich in einer Opposition gegenüber und schwarz muss ziehen. \\
Weiß kann mitziehen. Wenn schwarz sich nah an einem Bauern befindet, kann dieser nicht weiter und muss den Bauern aufgeben. \\ Er bekommt zwar den Bauern auf der anderen Seite. Allerdings muss er dahin noch laufen. \\
\centering
\begin{tabular}[h]{c|c}
	\textbf{Weiß}  & \textbf{Schwarz} \\
	\hline
	Ke2 & Ke7 \\
	Ke3 & Ke6 \\
	Ke4 & Kd6 \\
	Kd4 & Ke6 \\
	Kc5 & $\cdot$\\
	\unterstreichen{1-0}
\end{tabular}



\section{Kapitel 16}
\white{8/8/8/4p1p1/8/5P2/6K1/3k4 w - - 0 1}\\
\links

Es darf keine Stellung entstehen wo der weiße König auf g3 und der schwarze auf e3 steht (und weiß dran ist). So eine Stellung verliert weiß immer.\\
Deswegen muss Weiß versuchen immer in Opposition zu gehen, damit dieser Fall nicht eintritt. \\
\centering
\begin{tabular}[h]{c|c}
	\textbf{Weiß}  & \textbf{Schwarz} \\
	\hline
	Kh1 & Ke1 \\
	Kg1 & Ke2 \\
	Kg2 & Ke3 \\
	Kg3 & Kd2 \\
	Kh2 & Kd3 \\
	Kh3 & $\cdot$ \\
	\unterstreichen{1-0}
\end{tabular}

\section{Kapitel 17}
\white{8/3k2p1/8/5P2/8/8/8/1K6 w - - 0 1}\\
\links

Hier ist es wichtig zu erkennen, wann ich den Hebel f6 spiele. \\
Wenn ich gleich f6 spiele, kommt der Bauer zwar auf die f-Linie, allerdings habe ich dann keine Möglichkeit mehr mit dem schwarzen König in Opposition zu gehen. Deswegen muss ich mit Kc2 anfangen.\\
Danach geht alles wie gewohnt.\\

\centering
\begin{tabular}[h]{c|c}
	\textbf{Weiß}  & \textbf{Schwarz} \\
	\hline
	Kc2 & Kd6 \\
	f6 & gxf6 \\
	Kd2 & Ke6 \\
	Ke2 & f5\\
	Kf3 & Ke5 \\
	Ke3 & $\cdot$ \\
	\unterstreichen{\Remis} & \unterstreichen{\Remis}
\end{tabular}



\section{Kapitel 18}
\black{8/5p2/8/6Pk/5P2/8/8/7K w - - 0 1}\\
\links

Wenn ich g6 spiele, muss entweder der König oder der Bauer nehmen. Das selbe gilt für das selbige f5. \\
Danach kann ich mithilfe von Kf1 die Opposition halten.\\

\centering
\begin{tabular}[h]{c|c}
	\textbf{Weiß}  & \textbf{Schwarz} \\
	\hline
	Kc2 & Kd6 \\
	f6 & gxf6 \\
	Kd2 & Ke6 \\
	Ke2 & f5 \\
	Kf3 & Ke5 \\
	Ke3 & $\cdot$ \\
	\unterstreichen{\Remis} & \unterstreichen{\Remis}
\end{tabular}


\section{Kapitel 19}
\white{8/6k1/3K1p2/4p1p1/4P1P1/5P2/8/8 b - - 0 1}\\
\links
Auch hier gilt, Oppositionen sind nice.\\
Also bleibt hier nur das Feld Kh6. Und dann einfach mitziehen. \\

\centering
\begin{tabular}[h]{c|c}
	\textbf{Weiß}  & \textbf{Schwarz} \\
	\hline
	$\cdots$ & Kh6\\
	f6 & gxf6 \\
	Ke6 & Kg6 \\
	Ke7 & Kg7 \\
	\unterstreichen{\Remis} & \unterstreichen{\Remis}
\end{tabular}

\chapter{Training 29.12.2020}
 Die folgenden Aufgaben beziehen sich auf die zweite Lichess Studie \cite{Study2}

\vspace{5cm}
\hspace{2cm}

\begin{figure}[h]
	\centering
	\includegraphics[height=6cm]{question.png}
	\caption{Question}
	\label{fig:Question}
\end{figure}

% Die einzelnen Kapitel
\cleardoublepage % Seitenumbruch erzwingen vor Änderung des Nummerierungsstils
\pagenumbering{arabic} % Nummerierung der Seiten ab hier: 1, 2, 3, 4...


\appendix
\part*{Anhang}




% Literaturverzeichnis einbinden
%% ++++++++++++++++++++++++++++++++++++++++++++++++++++++++++++
%% Anhang: Literaturverzeichnis
%% ++++++++++++++++++++++++++++++++++++++++++++++++++++++++++++


% Mit dem Befehl \nocite werden auch nicht im Text zitierte
% aus der Literaturdatenbank mit in das Literaturverzeichnis aufgenommen.
% Ein "\nocite{*}" übernimmt ungeprüft die komplette Datenbank.
%\nocite{*}

\cleardoublepage
\ihead[]{Literaturverzeichnis}
\bibliographystyle{alphadin}
\bibliography{literatur} % "literatur.bib" ist hier die einzige Literaturdatenbank.

% Alternativ: Mehrere Datenbanken verwenden, falls eine
% oder mehrere umfangreiche Sammlungen exisitieren:
%\bibliography{literatur_buecher,literatur_weblinks}


% Abbildungsverzeichnis einbinden
%% ++++++++++++++++++++++++++++++++++++++++++++++++++++++++++++
%% Anhang: Abbildungsverzeichnis
%% ++++++++++++++++++++++++++++++++++++++++++++++++++++++++++++


% Keine Änderungen vornehmen!
\cleardoublepage
\ihead[]{Abbildungsverzeichnis}
\listoffigures


% Tabellenverzeichnis einbinden
%%% ++++++++++++++++++++++++++++++++++++++++++++++++++++++++++++
%% Anhang: Tabellenverzeichnis
%% ++++++++++++++++++++++++++++++++++++++++++++++++++++++++++++

% Hier keine weiteren Änderungen vornehmen
\cleardoublepage
\ihead[]{Tabellenverzeichnis}
\listoftables



% Abkürzungsverzeichnis einbinden
%% ++++++++++++++++++++++++++++++++++++++++++++++++++++++++++++
%% Anhang: Abkürzungsverzeichnis
%% ++++++++++++++++++++++++++++++++++++++++++++++++++++++++++++


% Hier keine weiteren Änderungen vornehmen
\cleardoublepage
\ihead[]{Abkürzungsverzeichnis und Formelzeichen}
\printnomenclature





% Abschlusserklärung
%%% ++++++++++++++++++++++++++++++++++++++++++++++++++++++++++++
%% Vorwort: Abschlusserklärung
%% ++++++++++++++++++++++++++++++++++++++++++++++++++++++++++++


\chapter*{Erklärung}
\addcontentsline{toc}{chapter}{Erklärung}
\ihead[]{Erklärung}

Die vorliegende Arbeit habe ich selbstständig ohne Benutzung anderer als der
angegebenen Quellen angefertigt. Alle Stellen, die wörtlich oder sinngemäß
aus veröffentlichten Quellen entnommen wurden, sind als solche
kenntlich gemacht. Die Arbeit ist in gleicher oder ähnlicher Form oder
auszugsweise im Rahmen einer oder anderer Prüfungen noch nicht vorgelegt
worden.
\\[2cm]
Ilmenau, den \today \hfill \namedesautors


\end{document}

